% META: {"id":"TEXP-MATRICES-MARKOV-RE-C5","chapitre":"TEXP-MATRICES-MARKOV","type_objet":"remediation","capacites_codes":["C5"],"status":"approved","origin":"generated","status_history":["generated","audited","accepted","approved"],"release_acceptance":"RELEASE_OWNER_BATCH_ACCEPTANCE","acceptance_closure_digest":"sha256:9b3ccf9a81c5520fbb7e03b7b2d3e7bf2057a02834b6d908bf3be5f49f3b553f","acceptance_receipt_digest":"sha256:4fad86f0282e0fa4e0419cca1ad6379ae7af5a280c04a4a90880f90a1c044242"}

%---------------------------------------------------------------
% FICHE DE REMEDIATION — C5 : pi_0 P^n (ligne x matrice), pas P^n pi_0
%---------------------------------------------------------------

\begin{exercice}{TEXP-MAT-RE-C5-EX1}{1}{10}
Un élève calcule la distribution après $2$ transitions en écrivant $\pi_0$ comme une matrice \textbf{colonne} et en calculant $P^2\pi_0$ (au lieu de $\pi_0P^2$ avec $\pi_0$ en matrice ligne).

On reprend la chaîne de Markov du streaming : $P=\begin{pmatrix}0{,}9&0{,}1\\0{,}3&0{,}7\end{pmatrix}$, $\pi_0=\begin{pmatrix}1&0\end{pmatrix}$.
\begin{enumerate}
  \item Calculer la distribution correcte $\pi_2=\pi_0P^2$ (avec $\pi_0$ en matrice ligne, multipliée \textbf{à gauche} de $P^2$).
  \item Calculer $P^2\begin{pmatrix}1\\0\end{pmatrix}$ (en transformant, comme l'élève, $\pi_0$ en matrice colonne multipliée \textbf{à droite}).
  \item Comparer les deux résultats. Que peut-on en conclure sur la méthode de l'élève ?
\end{enumerate}
\end{exercice}

% BEGIN-VERIFY
% from sympy import Matrix, Rational
% P = Matrix([[Rational(9,10), Rational(1,10)], [Rational(3,10), Rational(7,10)]])
% pi0_row = Matrix([[1,0]])
% pi0_col = Matrix([[1],[0]])
% correct = pi0_row * P**2
% wrong = P**2 * pi0_col
% assert correct == Matrix([[Rational(21,25), Rational(4,25)]])
% assert wrong == Matrix([[Rational(21,25)],[Rational(12,25)]])
% assert correct[0,1] != wrong[1,0]
% END-VERIFY

\begin{corrige}{TEXP-MAT-RE-C5-EX1}
\textbf{1.} $\pi_2=\pi_0P^2=\begin{pmatrix}1&0\end{pmatrix}\begin{pmatrix}0{,}84&0{,}16\\0{,}48&0{,}52\end{pmatrix}=\begin{pmatrix}0{,}84&0{,}16\end{pmatrix}$ : la probabilité correcte d'être non-abonné après $2$ mois est $0{,}16$.

\textbf{2.} $P^2\begin{pmatrix}1\\0\end{pmatrix}=\begin{pmatrix}0{,}84\\0{,}48\end{pmatrix}$ : la seconde coordonnée obtenue par ce calcul est $0{,}48$.

\textbf{3.} Les deux résultats diffèrent ($0{,}16$ contre $0{,}48$) : la méthode de l'élève est \textbf{fausse}. En transformant $\pi_0$ en matrice colonne multipliée à droite de $P^2$, on calcule en réalité la \textbf{première colonne} de $P^2$, qui n'a pas le même sens probabiliste que $\pi_0P^2$. La distribution $\pi_0$ doit toujours rester une matrice \textbf{ligne}, multipliée \textbf{à gauche} de $P^n$.
\end{corrige}
